Three parts, each machine-checked. \emph{The algebra of $S$.} For $B$ in the plane pair: $\sigma_r\,\sigma_r\,\sigma_r = \sigma_r$ since $\sigma_r^2 = 1$, and $\sigma_r (I\sigma_r)\sigma_r = I\sigma_r^3 = I\sigma_r$ since $I$ commutes with every even element (\claimref{std-algebra.pseudoscalar-commutation}); so $-\tfrac12(B + 3\sigma_r B\sigma_r) = -2B$ there. For each of the four mixed planes, $B$ shares exactly one index with the radial plane, hence anticommutes with $\sigma_r$ (\claimref{std-algebra.orthogonal-anticommute}): $\sigma_r B \sigma_r = -B\sigma_r^2 = -B$, giving $+B$. This is $1 - 3\Pi$; then $S^2 = 1 + 3\Pi$ (idempotency of $\Pi$) gives $S^2 + S = 2$, the spectrum is read off, and the trace is $2\cdot(-2) + 4\cdot 1 = 0$. \emph{The eigenvalues.} On the radial plane the eigenvalue comes from the gauge field alone: Proposition \claimref{singularity.solution-where-divergence-lies}(1) computed, in two lines and with the commutator term vanishing identically,
\[
\mathcal{R}(\partial_t \wedge \partial_r) \;=\; -\,\partial_r\Bigl(\frac{M}{r^2}\Bigr)\,\sigma_r \;=\; \frac{2M}{r^3}\,\sigma_r .
\]
The mixed and angular planes are computed in Appendix \ref{app:curvature} by the Cartan structure equations in the static chart, giving $-M/r^3$ on the four mixed planes and the doubled value again on the angular dual; curvature is a pointwise invariant, and the two pictures agree on the invertible region by Fact \claimref{singularity.imported-equivalence-einstein-equations}. The eigenvalue ratio is $-2 : 1$, which is the ratio of $S$'s spectrum. \emph{Assembly.} An operator diagonal on the six planes with eigenvalues proportional to $(-2,-2,1,1,1,1)$ is $\pm(M/r^3)(1 - 3\Pi)$, and the sign is fixed by the radial plane: $\mathcal{R}(\sigma_r) = +2M/r^{3}\,\sigma_r$ against $S(\sigma_r) = -2\sigma_r$ gives the minus sign, and the mixed planes, $-M/r^{3}$ against $S = +1$, confirm it. The $\hb$-invariance display is one line of anticommutation ($\er \wedge \er = 0$). The script checks the conjugation table, the operator identities, the spectrum and trace, the eigenvalue ratio against the appendix's rational-function values, the sign of the identity on the radial and mixed planes, and the invariance display, blade-exactly.
